% id: 38-23
% book: 베이직쎈 미적분1 · 02 함수의 연속
% section: 개념 18 사잇값 정리
% figure: none
% answer: 
% answer_source: 
% variant_level: 0 (원본 전사)
% 이 파일은 build.mjs 가 items.json 에서 만든다. 고칠 때는 items.json 을 고친다.
\smash{\raisebox{1.5mm}{\dmkichul{베이직쎈 미적분1 38쪽 23번}}}
다음은 함수 $f(x)=x^2-1$에 대하여 $f(c)=1$인 $c$가 열린구간 $(1{,}\mkern3mu\mathopen{}2)$에 적어도 하나 존재함을 증명한 것이다. \blank\ 안에 알맞은 것을 써넣으시오.\exprbox{\begin{minipage}{0.96\linewidth}{\small\bfseries 증명}\par\smallskip\raggedright 함수 $f(x)=x^2-1$은 닫힌구간 $[1{,}\mkern3mu\mathopen{}2]$에서 \blank[2.4em]이다. 또 $f(1)=\blank$, $f(2)=\blank$에서\par\smallskip\dispeq{$f(1)\ne f(2)$}\par\smallskip 이고 $f(1)<\blank<f(2)$이므로 사잇값 정리에 의하여 $f(c)=1$인 $c$가 열린구간 $(1{,}\mkern3mu\mathopen{}2)$에 적어도 하나 존재한다.\end{minipage}}
